\section{\faUsers\ Work Experience}

\datedsubsection{WeRide, Perception R\&D Engineer}{2024.06 - Present}

\paragraph{AEB Obstacle Perception Fusion}
\begin{itemize}
      \item Participated in the development of AEB functionality for Chery \textit{single Orin Y} model and GAC \textit{G300} model,
            designed multi-sensor fusion algorithms.
      \item ASIL-B requirements mandate fault tolerance through input redundancy.
            Designed multi-source fusion algorithm integrating vision end-to-end model,
            vision AEB model, and Radar detection (\textit{G300}).
      \item Fused multi-source detections based on Kalman Filter.
            Decoupled sensor-specific characteristics through measurement noise modeling.
            Aligned timestamps through temporal prediction.
      \item Orin Y model delivery achieved 100\% UT coverage,
            5k hours simulation with zero false positives, and passed C-NCAP assessment.
\end{itemize}

\paragraph{LiDAR Freespace Post-processing}
\begin{itemize}
      \item High-reflectivity boards (road signs, traffic cones, etc.)
            cause strong false alarms in LiDAR perception under rainy conditions.
            Designed rule-based post-processing to remove FP.
            Fit the plane where high-reflectivity board point clouds are located through PCA,
            remove freespace false positives caused by high-reflectivity boards.
            Compensate with vision model outputs.
      \item Implemented freespace memory module based on occupancy grid tracking.
            Completed freespace behind the ego vehicle through blind spot traversal determination strategy,
            providing stable and predictable results compared to Log-Odds probability accumulation.
            Supported complex maneuvers such as three-point turns and filtered sporadic false detections.
      \item Designed conversion algorithm from freespace point cloud to linestring for in-vehicle HMI rendering.
            Based on stable results provided by freespace memory module,
            quantize angles with adaptive density, compute convex hull of sampled results,
            then enrich details by adding freespace points between convex hull vertices.
            Enable in-vehicle HMI to render freespace contour lines while maintaining accuracy.
\end{itemize}

\paragraph{NeuralMap Lane Line Post-processing}
\begin{itemize}
      \item NeuralMap is the lane detection module of WeRide compass mapless architecture.
            Used in all L2 mapless vehicle models in collaboration with Chery, Bosch, etc.
      \item Developed and adjusted post-processing features
            based on downstream feedback and requirements.
            Assigned attributes such as merge/diverge, lane marking type, etc.
      \item Recognized crosswalks and smoothed results based on Kalman Filter,
            providing stable results for HMI rendering.
      \item Developed bag-based labeling tool,
            implemented automatic recognition and labeling of ego vehicle behaviors
            (lane centering/off-centering, lane change, stop line stationary/start/crossing)
            and interactive vehicle behaviors (stationary, acceleration/deceleration, cut-in, bypass).
\end{itemize}



\datedsubsection{TuSimple, R\&D Engineer - HD Map (full time \& intern)}{2022.07 - 2024.03}

\paragraph{Maintenance and Development of HD Map SDK TSMap}
\begin{itemize}
      \item TSMap is TuSimple's HD map C++ SDK,
            which also provides Python interface through pybind11.
      \item Developed Perception Drivable Map.
            Provides interface for perception to quickly obtain boolean matrix representing drivable area.
      \item Developed Planning Drivable Map.
            Compute signed distance field to freespace boundaries.
            Utilize spatial locality to asynchronously compute results meeting real-time requirements.
      \item Performed optimizations based on downstream use cases.
            For example, provided batch curvature query interface, optimized algorithmic complexity for batch queries.
\end{itemize}

\paragraph{Map Production Pipeline}
\begin{itemize}
      \item Participated in the development of new map production pipeline with microservice architecture
            for large-scale map production.
      \item Developed lane-level map construction related services.
            Pipeline inputs lane line vector data,
            outputs HD map basedata containing map elements such as lanes, intersections,
            lane-level topology, physical element relationships, etc.
            Integrated with observability platform for rapid issue discovery and root cause analysis.
      \item Maintained legacy map compilation tool and developed new features.
            Extended map element types according to perception and planning requirements,
            supporting representation of double-line boundaries, non-lane drivable areas, etc.,
            while ensuring forward and backward compatibility.
\end{itemize}
